\section{What this means, and what remains open}\label{sec:limits}

\subsection{What this means}

We set out to learn what a fast decision model buys an agent harness. The answer has three parts. Decisions are fast
and cheap, so speed is never the obstacle. On labelled decisions the best-value model is Jev, but on real read
decisions no model was safe enough to act alone. And where the money is, the model a whole session runs on, the
measured saving came from price arithmetic and a deterministic rule, not from the decision model: routing to Sonnet
at medium effort saves on Fable (\SHOneRatio\X, \cref{sec:plateau}), never routing is right on Opus, and the rule
matched the model at the session start.

\paragraph{The configuration answer} Decide once per session, keep the price gate on, use the rule R*, set effort to
medium on Sonnet and Fable but not Opus, and keep review and explain work on the host (\cref{sec:recommend}).

\subsection{What remains open}


The plan for this program (\path{docs/design/v3/PLAN.md}) has five sessions studies; only the first is in this
report. Each of the others has a written design there:
\begin{itemize}
  \item \textbf{S2, long sessions:} does decide-once still hold when a session runs far beyond \MaxTurns\ turns?
  \item \textbf{S3, delegation in situ:} routing the sub-tasks an agent delegates, inside real sessions.
  \item \textbf{S4, other harnesses:} the measured saving of \code{decide} and \code{launch} in Claude Code, Codex
    and Copilot CLI.
  \item \textbf{S5, SWE-bench at full scale:} two arms on the full benchmark.
  \item \textbf{Replicating the Fable freeze choice in S2} (it was selected on holdout), and anything on hosts other
    than Fable, Opus and Sonnet.
\end{itemize}

\paragraph{Honest limits.} Everything here ran on one provider, at one price table: costs are list prices applied
to measured tokens, tools-normalized as described in \cref{sec:plateau}, not invoices, and the Opus conclusion depends on its cache-read price. The sessions are scripted,
at most \MaxTurns\ turns, on public code; the counterfactual reuses main-v1's sessions and its cheap outcomes for a
few scenario-reps are estimates. The telemetry is one owner's workload. Decision quality on real reads remains
unidentified beyond the \TrN\ labelled cases. The judge's own cost is small but was estimated, not billed. S1 had
operational deviations (\cref{sec:plateau}): a raised budget cap, a stop rule not followed to the letter, and many
infrastructure reruns.

\paragraph{Where a decision model might still matter} The session-start choice is made once, with little information,
and a rule already captures it. A decision model may still earn its keep where the right answer changes during a
session: escalating a stuck session to the host, choosing which sub-task to delegate, or stopping a loop. That is a
hypothesis for S2--S4, not a finding.

\begin{keyidea}[title={The takeaway}]
At the start of a session, price arithmetic and a one-line rule captured all of the measured saving; a fast decision
model will have to prove its worth on decisions that change during the session.
\end{keyidea}
